%------------------------------------------------------------------------------%
\section{Identidades Trigonométricas}

Algumas relações conhecidas como \emph{Identidades Trigonométricas} serão
fundamentais no curso

%------------------------------------------------------------------------------%
\emph{Seno da soma}
\begin{equation}
  \label{eq:soma_seno}
  \sin(a + b) = \sin a \cos b + \sin b \cos a.
\end{equation}

%------------------------------------------------------------------------------%
\emph{Cosseno da soma}
\begin{equation}
  \label{eq:soma_cosseno}
  \cos(a + b) = \cos a\cos b - \sin b \sin a
\end{equation}

%------------------------------------------------------------------------------%
\OnlineBox{https://www.youtube.com/watch?v=MlCZU4zNymo}
{
  Para visualizar as demonstrações das identidades~\eqref{eq:soma_seno}
  e~\eqref{eq:soma_cosseno}, assista ao vídeo disponível neste link.
}
%------------------------------------------------------------------------------%

%------------------------------------------------------------------------------%
\emph{Tangente da soma}
\begin{equation}
  \label{eq:soma_tangente}
  \tan(a + b) = \frac{\tan a + \tan b}{1 - \tan a \tan b}.
\end{equation}

%------------------------------------------------------------------------------%
\begin{proof}
  Para demonstrar esse item, usamos a definição de função tangente e
  as expressões~\eqref{eq:soma_seno} e~\eqref{eq:soma_cosseno}. Assim,
  \[
    \tan(a+b)
    = \frac{\sin(a+b)}{\cos(a+b)}
    = \frac{\sin a\cos b+\sin b\cos a}{\cos a\cos b-\sin b\sin a}
  \]
  Dividindo o numerador e o denominador por $\cos a\cos b$, desde que
  seja não nulo, temos
  \[
    \tan(a+b)
    = \frac{\;\dfrac{\sin a\cos b+\sin b\cos a}{\cos a\cos b}\;}
           {\;\dfrac{\cos a\cos b-\sin b\sin a}{\cos a\cos b}\;}
    = \frac{\tan a+\tan b}{1-\tan a\tan b}
  \]
  Portanto, a fórmula para a tangente da soma é
  \[
    \tan(a+b)=\frac{\tan a+\tan b}{1-\tan a\tan b}
  \]
\end{proof}
%------------------------------------------------------------------------------%

%------------------------------------------------------------------------------%
\emph{Seno da diferença}
\begin{equation}
  \label{eq:diferenca_seno}
  \sin(a-b)=\sin a\cos b-\sin b\cos a.
\end{equation}

%------------------------------------------------------------------------------%
\emph{Cosseno da diferença}
\begin{equation}
  \label{eq:diferenca_cosseno}
  \cos(a-b)=\cos a\cos b+\sin b\sin a.
\end{equation}

As identidades~\eqref{eq:diferenca_seno} e~\eqref{eq:diferenca_cosseno}
podem ser demonstrados trocando $b$ por $-b$ nas
expressões~\eqref{eq:soma_seno} e~\eqref{eq:soma_cosseno},
respectivamente, e usando que a função seno é ímpar e a função cosseno é par.

%------------------------------------------------------------------------------%
\emph{Arco duplo: Seno}
\begin{equation}
  \label{eq:arco_duplo_seno}
  \sin(2a) = 2\sin(a)\cos(a).
\end{equation}

Essa fórmula é muito útil na trigonometria, pois permite calcular o
valor do seno de um ângulo duplo em termos dos valores de seno e
cosseno do ângulo original.

%------------------------------------------------------------------------------%
\begin{proof}
  Para demonstrá-la, vamos partir da fórmula da adição para o
  seno~\eqref{eq:soma_seno} , trocando $b$ por $a$. Logo,
  \begin{align*}
    \sin(2a)
    & = \sin(a + a)                      \\
    & = \sin(a)\cos(a) + \cos(a)\sin(a)  \\
    & = 2\sin(a)\cos(a)
  \end{align*}
  Portanto, a fórmula para o seno do arco duplo é
  \[
    \sin(2a) = 2\sin(a)\cos(a)
  \]
\end{proof}
%------------------------------------------------------------------------------%

%------------------------------------------------------------------------------%
\emph{Arco duplo: Cosseno}
\begin{equation}
  \label{eq:arco_duplo_cosseno}
  \cos(2a)=\cos^2(a)-\sin^2(a)
\end{equation}

%------------------------------------------------------------------------------%
\begin{proof}
    Vamos partir da fórmula da adição para o cosseno~\eqref{eq:soma_cosseno}
    \[
      \cos(\alpha + \beta)
      = \cos(\alpha)\cos(\beta) - \sin(\alpha)\sin(\beta)
    \]
    Substituindo $\alpha = a$ e $\beta = a$, temos
    \begin{align*}
      \cos(2a)
      = \cos(a + a)                     \\
      = \cos(a)\cos(a) - \sin(a)\sin(a) \\
      = \cos^2(a)-\sin^2(a).
    \end{align*}
    Portanto, a fórmula para o cosseno do arco duplo é
    \[
      \cos(2a) = cos^2(a)-\sin^2(a)
    \]
\end{proof}
%------------------------------------------------------------------------------%

Note que isolando $\sin^2(x)$ na identidade trigonométrica
\[
  \sin^2(a) + \cos^2(a) = 1,
\]
podemos reescrever a expressão~\eqref{eq:arco_duplo_cosseno} como
\[
  \cos(2a) = \cos^2(a) - \left(1 - \cos^2(a)\right)
\]
Simplificando, obtemos outra expressão para $\cos(2a)$, que depende
apenas de cosseno
\begin{equation}
  \label{eq:arco_duplo_cosseno_1}
  \cos(2a) = 2\cos^2(a) - 1
\end{equation}
Se isolarmos $\cos^2(x)$ na identidade trigonométrica
$\sin^2(a) + \cos^2(a) = 1$,
podemos reescrever a expressão \eqref{eq:arco_duplo_cosseno} como
\begin{equation}
  \label{eq:arco_duplo_cosseno_2}
  \cos(2a) = 1-2\sin^2(a)
\end{equation}
Essas expressões serão essenciais para demonstrar as fórmulas de arco
metade que seguem.

%------------------------------------------------------------------------------%
\emph{Arco metade: Cosseno}
\begin{equation}
  \label{eq:arco_metade_cosseno}
  \cos^2(x) = \frac{1}{2}\big(1+\cos(2x)\big).
\end{equation}

%------------------------------------------------------------------------------%
\begin{proof}
  Começamos com a fórmula~\eqref{eq:arco_duplo_cosseno_1},
  \[
    \cos(2a) = 2\cos^2(a) - 1
  \]
  e substituímos $a$ por $\sfrac{x}{2}$ para obter
  \[
    \cos(x) = 2\cos^2\left(\frac{x}{2}\right) - 1
  \]
  Isolando \(\cos^2\left(\frac{x}{2}\right)\), temos
  \[
    \cos^2\left(\frac{x}{2}\right) = \frac{1}{2}(1 + \cos(x))
  \]
  Substituindo $\sfrac{x}{2}$ por $x$ e reescrevendo a equação em
  termos de $\cos(2x)$, temos
  \[
    \cos^2(x) = \frac{1}{2}(1 + \cos(2x))
  \]
  Portanto,
  \[
    \cos^2(x)= \frac{1}{2}(1+\cos(2x))
  \]
\end{proof}
%------------------------------------------------------------------------------%

%------------------------------------------------------------------------------%
\emph{Arco metade: Seno}
\begin{equation}
  \label{eq:arco_metade_sseno}
  \sin^2(x)= \frac{1}{2}\big(1-\cos(2x)\big).
\end{equation}

%------------------------------------------------------------------------------%
\begin{proof}
  Para comprovar esse item, é possível utilizar a mesma abordagem
  utilizada anteriormente, porém aplicando agora a
  fórmula~\eqref{eq:arco_duplo_cosseno_2}.
  Dessa forma, será possível demonstrar de maneira
  equivalente a veracidade desse item.
\end{proof}
%------------------------------------------------------------------------------%
